\documentclass[preprint,12pt]{elsarticle}
% \documentclass{article}


\usepackage{amsmath}
\usepackage{amssymb}
\usepackage{graphicx}
\usepackage{caption}
\usepackage{array}
\usepackage{epstopdf}
\usepackage{hyperref}
\usepackage{booktabs}
\usepackage{geometry}
\usepackage{subcaption}
\usepackage{threeparttable}


% Define new highlight colors
\usepackage{color,soul}
\usepackage{xcolor}
\sethlcolor{yellow}

\colorlet{lightcyan}{cyan!30}    % 30% cyan, 70% white
\colorlet{lightgreen}{lime!30}   % 30% lime, 70% white
\colorlet{lightorange}{orange!35} % 35% orange, 65% white

% \colorlet{blue}{black}   % Clean (non-highlighted) version

\newcommand{\hlcyan}[1]{{\sethlcolor{lightcyan}\hl{#1}}}
\newcommand{\hlgreen}[1]{{\sethlcolor{lightgreen}\hl{#1}}}
\newcommand{\hlorange}[1]{{\sethlcolor{lightorange}\hl{#1}}}


% \journal{Energy Strategy Reviews}


\begin{document}

\begin{frontmatter}

\title{PyPSA-Libya: An Open-Source Energy System Model for National Decarbonization and Hydrogen Export Pathways}

\author[a]{Mohamed Abuella}
\ead{mhdabuella@gmail.com}

% \affiliation[a]{organization={Independent Researcher}, city={Newcastle upon Tyne}, country={UK}}

\begin{abstract}
Libya possesses significant renewable energy potential that remains largely untapped. This paper presents PyPSA-Libya, a national-scale model built on the PyPSA-Earth framework, designed to explore optimal transition pathways for the Libyan power sector. We present two primary studies: (1) an electricity-only capacity expansion analysis comparing a zero-emission constraint against an unconstrained baseline at 36-cluster resolution, and (2) a sector-coupled hydrogen export study evaluating the levelized cost of exporting 70 TWh/year of hydrogen to Europe under three technological scenarios. Results indicate that while solar PV dominates the capacity mix due to favorable 2030 cost assumptions, the integration of hydrogen production significantly alters the grid's storage requirements. The model, data, and configuration files are released under an open-source license to support transparent energy planning in Libya.
\end{abstract}

\begin{keyword}
Energy system modelling \sep Libya \sep Green Hydrogen \sep Decarbonization \sep PyPSA-Earth
\end{keyword}

\end{frontmatter}


\section{Introduction}
\label{sec:introduction}

The global imperative to mitigate climate change has accelerated the transition from carbon-intensive fossil fuels to renewable energy sources, in line with the objectives of the Paris Agreement \cite{almaktar2021prospects, alanazi2025renewable}. Petro-state economies face a distinct dual challenge during this transition: they must decarbonize their domestic power sectors while simultaneously diversifying their economies away from oil and gas dependence \cite{almaktar2021prospects, eljrushi1990solar}. Libya represents a critical case study of this paradigm. The country's political economy relies almost entirely on conventional fossil fuels, with the oil sector contributing approximately 95\% of export earnings and 65\% of Gross Domestic Product (GDP) \cite{almaktar2021prospects}. Over the past decade, political instability and administrative division have severely affected Libya's energy sector. Armed conflicts and unilateral shutdowns of major oil fields (such as those in early 2020 causing financial losses exceeding \$9 billion) have repeatedly undermined the reliability of fuel supplies for electricity generation \cite{almaktar2021prospects}.

Concurrently, Libya's domestic power sector, monopolized by the state-owned General Electricity Company of Libya (GECOL), is facing a severe operational crisis \cite{almaktar2021prospects}. The power grid relies completely on fossil fuels, with natural gas driving approximately 63\% of generation and oil-fired plants accounting for 29\% \cite{almaktar2021prospects}. Rapid growth in energy demand, driven by rising temperatures and population growth, has led to a persistent peak-load electricity deficit of approximately 1,200~MW, resulting in frequent and widespread blackouts \cite{almaktar2021prospects}. Furthermore, Libya's environmental footprint is substantial; in 2018, its per capita CO2 emissions amounted to 8.73 tonnes, surpassing the Chinese national average of 7.95 tonnes \cite{almaktar2021prospects}. Although the Renewable Energy Authority of Libya (REAoL) set an ambitious target in 2010 to integrate a 10\% renewable share (equivalent to 2,219~MW) into the national mix by 2025, this initiative was completely derailed by the civil conflict that erupted in 2011 \cite{almaktar2021prospects}.

Despite these challenges, Libya is exceptionally well-positioned to pivot toward a sustainable energy economy. The country possesses some of the highest solar and wind resource potentials in the world. The average global horizontal irradiation (GHI) over the Libyan territory yields an average potential solar photovoltaic (PV) productivity of 1.9~MWh/kW/year (with Al Kufra in the southeast reaching an optimal annual productivity of 2,100~kWh/kWp), while onshore wind potentials reach up to 500~W/m$^2$ along its southernmost borders, coastal Darna, and the Western Mountain region \cite{almaktar2021prospects}. Recognizing these vast resources, pioneering research in the late 1980s first proposed a solar-hydrogen energy system for Libya, emphasizing that converting solar power to hydrogen would not only extend the lifespan of fluid fossil fuel reserves as raw chemical feedstocks but also position Libya as a leading exporter of clean hydrogen to Europe in lieu of oil and gas \cite{eljrushi1990solar}. 

To evaluate the modern technical and economic feasibility of this transition, we introduce \texttt{PyPSA-Libya}, a high-resolution, open-source energy system model of Libya's power and hydrogen network. Built as a country-specific extension of the global \texttt{PyPSA-Earth} framework \cite{abdelkhalek2025pypsa}, our model is divided into two primary, complementary studies designed to address both domestic and international planning questions:
\begin{enumerate}
    \item {Electricity-only capacity expansion (Part 1):} This study models Libya's power sector in isolation at a high spatial resolution of 36 clustered nodes. It compares a zero-CO2 emission constraint (Case 1) against an unconstrained fossil-heavy baseline (Case 2) under 2030 technology-cost assumptions, capturing transmission bottlenecks and identifying the optimal mix of solar, wind, and storage needed for complete domestic grid decarbonization.
    \item {Sector-coupled hydrogen export (Part 2):} This study models the spatial and economic implications of exporting 70~TWh/year of renewable or low-carbon hydrogen from Libya to European markets via maritime shipping \cite{abdelkhalek2025pypsa}. Formulated at a 12-cluster resolution to maintain computational tractability, we evaluate three technology scenarios representing strictly green electrolysis (Scenario A), electrolysis coupled with carbon-captured steam methane reforming (SMR CC) representing blue hydrogen (Scenario B), and an unconstrained scenario representing grey, blue, and green pathways (Scenario C) .
\end{enumerate}

To ensure robustness, our model incorporates a \texttt{connectivity\_upgrade} workflow to resolve topological gaps and isolated buses inherent in OpenStreetMap (OSM) grid data for Libya. We report the Levelized Cost of Electricity (LCOE) and apply a paired-run differencing methodology to isolate the Net Levelized Cost of Hydrogen (Net LCOH) directly attributable to the export obligation. The model code, parameters, and results are released under an open-source license to establish a transparent, reproducible benchmark for national decarbonization and cross-border energy partnerships.

% =====================================================================

\section{Related work}
\label{sec:related_work}

The design of the \texttt{PyPSA-Libya} model builds upon three main pillars: (i) localized studies on Libyan renewable and hydrogen potentials, (ii) advancements in open-source spatial-temporal modeling, and (iii) global hydrogen trade pathways and underground storage economics.

\subsection{Renewable potential and hydrogen in Libya}
Early national energy assessments by Eljrushi and Veziroglu \cite{eljrushi1990solar} mathematically simulated the dynamic interaction between population growth and fossil fuel depletion, demonstrating that Libya would face severe energy supply deficits without a strategic transition toward solar-hydrogen production. Almaktar and Shaaban \cite{almaktar2021prospects} modernized these estimates through GIS and meteorological evaluations, categorizing Libya's solar, wind, biomass, and geothermal capacities as essential non-rivalry energy alternatives. 

While macro-level studies established these resource potentials, localized techno-economic analyses remain rare. Rahil et al. \cite{rahil2018techno} evaluated off-peak forecourt electrolysis in Darnah under 2015 and 2030 cost horizons; however, their scope was restricted to isolated refueling stations without capturing nationwide transmission grid dynamics or spatial sector coupling. Our model bridges this gap by representing the full national transmission topology and co-optimizing domestic and export infrastructure.

\subsection{Open-source modeling and the PyPSA framework}
Proprietary energy planning models (e.g., TIMES, PLEXOS) often present commercial license barriers that limit transparency and collaboration in developing nations. To address this, the Python for Power System Analysis (PyPSA) library was introduced as a fully open-source framework for high-resolution capacity expansion planning. While \texttt{PyPSA-Eur} pioneered continental-scale transmission and storage optimization, country-specific forks such as \texttt{PyPSA-Spain} \cite{gallego2025pypsaspain} demonstrated the necessity of integrating detailed national datasets and nested border-price models to accurately reflect local grid conditions.

Expanding open-source modeling globally, Parzen et al. \cite{parzen2023pypsa} developed \texttt{PyPSA-Earth}, an automated global framework for extracting high-voltage network topologies from OpenStreetMap and mapping spatiotemporal weather potentials. Abdel-Khalek et al. \cite{abdelkhalek2025pypsa} extended this platform into a sector-coupled model incorporating multi-carrier energy vectors (electricity, hydrogen, gas, biomass, oil, and heat) across industry, transport, and residential sectors. Their application establishing Egypt's potential to supply 16\% of European hydrogen demand by 2030 forms the methodological benchmark for our sector-coupled export study.

\subsection{Global hydrogen trade and geological storage}
Cross-border hydrogen trade involves significant techno-economic and infrastructural uncertainties \cite{fakhreddine2025global}. In the Euro-Mediterranean context, TIAM-ECN modeling by Fattahi et al. \cite{fattahi2024opportunities} identified North Africa as the most cost-effective pipeline supplier of gaseous hydrogen to Europe, with liquefied shipping providing complimentary supply diversification. Alanazi et al. \cite{alanazi2025renewable} confirmed that international trade could meet over 31\% of global decarbonized hydrogen demand by 2050, highlighting North Africa-to-Europe as a key trade corridor. Nevertheless, high conversion and transport losses dictate that hydrogen should be prioritized for hard-to-abate sectors, long-duration storage, and heavy industry rather than direct space heating or light mobility \cite{johnson2025realistic}.

Managing the volatility of renewable supply under rigid export commitments presents a major operational challenge \cite{terlouw2022large}. Achhammer et al. \cite{achhammer2026impact} addressed this by integrating GIS-mapped geological salt caverns into \texttt{PyPSA-Earth} for Tunisia, showing that underground hydrogen storage (UHS) reduces levelized export costs by ~3.7\%, dampens price volatility, and safeguards domestic energy security. By integrating both maritime shipping logistics and underground storage limits, \texttt{PyPSA-Libya} delivers a comprehensive evaluation of how these spatial and geological assets optimize Libya's export resilience.

% =====================================================================

\section{Methodology}
\label{sec:methodology}

The \texttt{PyPSA-Libya} model is developed using the open-source \texttt{PyPSA-Earth} framework for large-scale energy system optimization. The model minimizes total annualized system costs—comprising capital expenditures (CAPEX) and variable operational costs (OPEX)—across defined temporal and spatial horizons.

\subsection{Input data and network topology}
The model utilizes OpenStreetMap (OSM) data for high-voltage transmission lines and substations. Preliminary diagnostic testing revealed significant data gaps in the Libyan OSM extract, specifically regarding isolated buses in the southern and desert regions (e.g., Fezzan and Al Kufra). To ensure a fully interconnected national grid backbone, synthetic line connections were added to guarantee that every node retains at least one physical connection to the transmission system.

Renewable resource potentials are derived from the ERA5 reanalysis weather dataset for the year 2013. The spatial resolution varies by study to balance detail with computational tractability:
\begin{enumerate}
    \item {Electricity-only study:} 36 clustered nodes, providing high-fidelity insight into internal transmission bottlenecks.
    \item {Sector-coupled study:} 12 clustered nodes, optimized to ensure the complex linear programming (LP) problem fits within 16~GB of memory.
\end{enumerate}

\subsection{Techno-economic assumptions}
All technology cost assumptions are calibrated to a 2030 planning horizon using the \texttt{PyPSA} technology-data repository. We assume a standard Weighted Average Cost of Capital (WACC) of 7.1\% (discount rate). The temporal resolution is resampled to 3-hour steps for the electricity study and 6-hour steps for the hydrogen study. This resolution is identified as the ``sweet spot'' for capturing solar diurnal ramps and battery cycling while preventing the numerical instability observed in 1-hour resolution runs for large-scale matrices.

\subsection{Scenario design}
The study is divided into two distinct modeling parts:

\subsubsection{Part 1: Electricity-only decarbonization}
To establish a decarbonization baseline, we compare two cases within the electricity-only framework (no sector coupling):
\begin{itemize}
    \item {Case 1 (Zero-CO2):} Imposes a hard \texttt{co2\_limit} of 0.0, prohibiting all fossil dispatch (OCGT, CCGT, oil) and forcing a 100\% renewable system (Solar, Wind, Battery, and $\text{H}_2$ storage).
    \item {Case 2 (Baseline):} Imposes a non-binding cap ($1.0 \times 10^9$~tonnes), allowing the model to freely select the cost-minimal generation mix relying heavily on natural gas.
\end{itemize}

\subsubsection{Part 2: Sector-coupled hydrogen export}
Following the methodology of \cite{abdelkhalek2025pypsa}, we model the export of 70~TWh/year of hydrogen via maritime shipping. The study is structured into two distinct policy and optimization phases to evaluate how carbon pricing alters technological choices:

\noindent{Phase 1: Baseline technology scenarios (No $\text{CO}_2$ tax)}
We first evaluate production economics under existing baseline market structures with zero penalty on carbon venting, isolating the performance of three distinct technological tracks:
\begin{itemize}
    \item {Scenario A (Green H$_2$):} Only water electrolysis is allowed, enforcing a strict annual renewable matching constraint.
    \item {Scenario B (Blue H$_2$):} Force-limits production to electrolysis and Steam Methane Reforming with Carbon Capture (SMR CC) under a rigid $\text{CO}_2$ cap ($1.0 \times 10^7$~tonnes).
    \item {Scenario C (Grey H$_2$):} An unconstrained baseline pathway allowing electrolysis, SMR CC, and unabated (grey) SMR under a non-binding emissions limit.
\end{itemize}

\noindent{Phase 2: Hybrid policy optimization (with $\text{CO}_2$ tax)}
We then introduce a uniform carbon tax applied directly to all net atmospheric emissions. Rather than running the previous scenarios in isolation, this phase removes all technological boundaries to form a single, integrated hybrid co-optimized scenario. In this unified run, the model is granted access to all three technology pathways simultaneously. The optimizer dynamically co-optimizes investments across green electrolysis, blue SMR CC, and grey SMR, automatically selecting the optimal hybrid technology mix that minimizes the total system cost under the new carbon tax regime.



\subsection{Solver configuration}
Numerical stability and computational scale are critical challenges in full-year, sector-coupled models. We utilize the HiGHS solver configured with the interior-point method (IPM). For complex runs, we disable the crossover phase and apply relaxed feasibility and optimality tolerances ($1 \times 10^{-3}$) to prevent numerical stall. The exact linear programming matrix dimensions, variable counts, and computational scaling metrics for Scenarios A, B, and C are detailed in~\ref{sec:appendix_problem_scale}.

\subsection{Data and cost assumptions}
We utilize ERA5 reanalysis weather data for the year 2013, resampled to 3-hour or 6-hour intervals. Technology costs are calibrated to the \textbf{2030 horizon}. A critical modification in \texttt{PyPSA-Libya} is the \texttt{connectivity\_upgrade} setting, which ensures that every node in the OSM-derived topology has at least one connection, mitigating the issue of isolated buses common in Libyan data.

% =====================================================================

\section{Results: electricity-only baseline}

\subsection{Capacity of the energy mix}

[Insert discussion here: Case 1 likely shows a massive buildout of Solar PV and Battery storage compared to Case 2's reliance on CCGT/OCGT.]
\begin{figure}[ht]
    \centering
    % Placeholder for visualize_libya_year.py output
    \includegraphics[width=0.75\textwidth]{Figures/Elec_LCO2/fig_solved_capacity_map.pdf}
    \caption{Optimal capacity mix for the 36-cluster electricity-only model (Zero-CO2 Case).}
    \label{fig:elec_capacity}
\end{figure}


\subsection{Transmission bottlenecks}
Utilizing the shadow prices (dual variables) of the transmission constraints, we identify key corridors requiring reinforcement to facilitate 100\% renewable penetration.

\begin{figure}[ht]
    \centering
    \includegraphics[width=0.8\textwidth]{Figures/Elec_LCO2/fig_transmission_line_bottlenecks_map.pdf}
    \caption{Map of transmission line loading and identified bottlenecks in the Libyan grid.}
    \label{fig:bottlenecks}
\end{figure}
\subsection{System cost breakdown}
The overall annualized system cost increases by 31.4\% when transitioning from the unconstrained baseline to the carbon-limited scenario, rising from 3.75 to 4.93 billion EUR/year. 

\begin{figure}[ht]
    \centering
    \includegraphics[width=0.70\textwidth]{Figures/Elec_LCO2/fig_cost_by_scenario.pdf}
    \caption{Annualized system cost breakdown by technology for the Limited $\text{CO}_2$ and Unlimited $\text{CO}_2$ electricity-only scenarios.}
    \label{fig:cost_breakdown}
\end{figure}

As shown in Figure~\ref{fig:cost_breakdown}, the primary driver behind this cost divergence is the capital expenditure required to firm up renewable generation under strict emission limits. In the Unlimited $\text{CO}_2$ case, fossil-fuel units are heavily dispatched to meet demand, incurring substantial operational and fuel costs (0.95 billion EUR/year OPEX) while minimizing the need for energy storage CAPEX (0.001 billion EUR/year). 

Conversely, enforcing emission constraints (Limited $\text{CO}_2$) shifts the system from a fuel-intensive operating model to a capital-intensive, solar-dominated paradigm:

\begin{itemize}
    \item {Solar PV capital expansion:} Investment in Solar PV capital expenditure (CAPEX) increases more than threefold, rising from 0.42 to 1.46 billion EUR/year to exploit Libya's high solar irradiance. Conversely, Onshore Wind CAPEX decreases from 0.62 to 0.40 billion EUR/year as Solar PV takes precedence.
    \item {Storage integration:} To bridge diurnal supply gaps without relying on high-emitting thermal dispatch, Battery Storage CAPEX surges from a negligible 0.001 billion EUR/year to 1.3138 billion EUR/year, accounting for 26.7\% of total system expenditure under the Limited $\text{CO}_2$ constraint.
    \item {Fossil displacement to reserve backup:} While CCGT and oil capital commitments (CAPEX) remain identical across both cases at 1.4972 and 0.2518 billion EUR/year, respectively, their operational role switches from baseload/mid-merit generation to low-utilization reserve backup. Consequently, system-wide variable operating and fuel costs (OPEX) plummet to negligible levels (0.0013 billion EUR/year).
\end{itemize}
\subsection{Financial analysis}
Enforcing zero $\text{CO}_2$ emissions requires a substantial upfront capital commitment to build out renewable energy generation and storage capacity. In contrast, the unconstrained, fossil-heavy system relies on significantly lower initial investment but incurs high ongoing operational and fuel expenses. 

To evaluate the economic trade-offs of full grid decarbonization in Libya, a financial feasibility comparison was conducted between the {Zero-$\text{CO}_2$ (Limited)} scenario and the unlimited baseline (unconstrained $\text{CO}_2$ emissions). Both configurations consider a 25-year operational horizon, a weighted average cost of capital (WACC) of 7\%, an annual demand of 53.65~TWh, and a benchmark electricity market price of \texteuro 75/MWh. A summary of the financial parameters, annual cash flows, payback dynamics, and $\text{CO}_2$ emissions balance is presented in Table~\ref{tab:financial_co2_comparison}.


\begin{table}[htbp]
\centering
\begin{threeparttable}
\caption{Financial parameters, economic performance, and annual $\text{CO}_2$ emissions across the Zero-$\text{CO}_2$ and Unconstrained scenarios.}
\label{tab:financial_co2_comparison}
\small
\begin{tabular}{l r r}
\toprule
{Metric} & {Zero-$\text{CO}_2$ (Limited)} & {Unlimited} \\
\midrule
\multicolumn{3}{l}{\textit{{A. Financial Parameters \& System Costs}}} \\
Overnight CAPEX (\texteuro B)\tnote{a} & 57.494 & 32.619 \\
Annualized CAPEX (\texteuro B/yr) & 4.934 & 2.799 \\
Annual OPEX (\texteuro M/yr) & 1.330 & 953.920 \\
Total Annual System Cost (\texteuro B/yr) & 3.162 & 1.980 \\
{LCOE (\texteuro/MWh)} & {58.94} & {36.90} \\
\midrule
\multicolumn{3}{l}{\textit{{B. Profitability \& Valuation (at \texteuro 75/MWh Target Tariff)}}} \\
Annual Projected Revenue (\texteuro B/yr)\tnote{b} & 4.024 & 4.024 \\
Annual Net Cash Flow (\texteuro B/yr) & 4.022 & 3.070 \\
Net Present Value (NPV @ 7\%)\tnote{c} (\texteuro B) & -10.617 & 3.157 \\
Internal Rate of Return (IRR) & 4.86\% & 8.05\% \\
Simple Payback Period (Years)\tnote{d} & 14.29 & 10.63 \\
Discounted Payback Period (Years)\tnote{e} & N/A & 20.12 \\
\midrule
\multicolumn{3}{l}{\textit{{C. Annual $\text{CO}_2$ Emissions Summary}}} \\
Total Annual Emissions ($\text{MtCO}_2$/yr) & 0.000 & 6.950 \\
System Carbon Intensity ($\text{g CO}_2$/kWh) & 0.00 & 0.13 \\
\bottomrule
\end{tabular}

\begin{tablenotes}
\footnotesize
\item[a] Overnight CAPEX represents the initial capital expenditure deployed at Year 0.
\item[b] Projected revenue is calculated using an assumed wholesale electricity price (target tariff) of \texteuro 75/MWh applied to the total annual electricity demand of 53.65~TWh.
\item[c] NPV is calculated over a 25-year project horizon ($T = 25$) at a 7.0\% discount rate. It represents the single cumulative present value of all cash flows evaluated at Year 0: $\text{NPV} = -\text{CAPEX} + \sum_{t=1}^{25} \frac{\text{Revenue} - \text{OPEX}}{(1 + r)^t}$, where Year 0 reflects upfront capital outlay and Years 1--25 reflect discounted operational net cash flows.
\item[d] Simple Payback Period represents the undiscounted recovery timeframe in years, computed as $\text{Simple Payback} = \frac{\text{CAPEX}}{\text{Annual Net Cash Flow}}$.
\item[e] Discounted Payback Period (DPP) accounts for the time value of money at discount rate $r = 7\%$, computed analytically as $\text{DPP} = \frac{-\ln\left(1 - \frac{\text{CAPEX} \cdot r}{\text{Annual Net Cash Flow}}\right)}{\ln(1 + r)}$. For Zero-$\text{CO}_2$, payback is undefined (N/A) as discounted annual inflows do not cover annual capital costs at $r = 7\%$.
\end{tablenotes}
\end{threeparttable}
\end{table}

\subsection{Economic Takeaways and Policy Implications}
The financial and emission findings highlight several key structural insights regarding the energy transition in Libya:

\begin{itemize}
    \item {Front-loaded CAPEX vs. OPEX trade-off:} Achieving a zero-emission grid inflates overnight capital expenditures to \texteuro 57.494~B (compared to \texteuro 32.619~B) due to extensive solar PV, wind, and battery storage investments. However, annual operational expenditures drop significantly to \texteuro 1.33~M (down from \texteuro 953.92~M) as fossil fuel purchasing is virtually eliminated.
    \item {The "Green premium" and tariff discrepancy:} At the baseline electricity market price of \texteuro 75/MWh, the unconstrained fossil system yields an LCOE of \texteuro 36.90/MWh (resulting in a positive NPV of \texteuro 3.157~B, an Internal Rate of Return (IRR) of 8.05\%, and a payback period of 10.63 years). Conversely, full decarbonization drives the LCOE to \texteuro 58.94/MWh (yielding an IRR of 4.86\% and extending the payback period to 14.29 years), establishing a green premium cost of \texteuro 22.04/MWh.
    \item {Policy and financing requirements:} Although the Zero-$\text{CO}_2$ grid generates higher net annual operational cash flow (\texteuro 4.022~B vs. \texteuro 3.070~B), the extended payback period (14.29 years) and negative NPV (\texteuro -10.617~B) demonstrate that it fails to fully recover its massive initial investment under standard commercial financing (7\% WACC). To make zero-carbon infrastructure bankable (NPV $\ge 0$), targeted policy interventions are required—such as securing long-term wholesale power purchase agreements (PPAs) above \texteuro 58.94/MWh, implementing carbon taxes on fossil fuels, or accessing concessional development financing with discount rates aligned with or below the project IRR of 4.86\%.
\end{itemize}

% =====================================================================

\section{Results: Sector-coupled hydrogen export} 
\label{sec:results_sector_coupled_h2_3secns} 
This section presents the results for a 70~TWh hydrogen export target across Scenarios A, B, and C, highlighting how alternative system configurations influence generation capacity, spatial investment patterns, and economic and environmental performance.

\subsection{Energy mix by carrier}
The cross-scenario comparison highlights how power supply dispatch transitions across technology carriers under varying system constraints.

Figure~\ref{fig:sect_coupl_capacity_all_scenarios} illustrates the generation mix by carrier across Scenarios A, B, and C, highlighting the differences in capacity deployment under the respective system constraints.

\begin{figure}[htbp]
    \centering
    \begin{subfigure}[b]{0.30\textwidth}
        \centering
        \includegraphics[width=\textwidth]{Figures/Sector_Couple_H2/fig_Scenario_A_physical_generation_by_carrier_filtered.pdf}
        \caption{Scenario A}
        \label{fig:sect_coupl_capacity_sen_a}
    \end{subfigure}
    \hfill
    \begin{subfigure}[b]{0.30\textwidth}
        \centering
        \includegraphics[width=\textwidth]{Figures/Sector_Couple_H2/fig_Scenario_B_physical_generation_by_carrier_filtered.pdf}
        \caption{Scenario B}
        \label{fig:sect_coupl_capacity_sen_b}
    \end{subfigure}
    \hfill
    \begin{subfigure}[b]{0.30\textwidth}
        \centering
        \includegraphics[width=\textwidth]{Figures/Sector_Couple_H2/fig_Scenario_C_physical_generation_by_carrier_filtered.pdf}
        \caption{Scenario C}
        \label{fig:sect_coupl_capacity_sen_c}
    \end{subfigure}
    \caption{Capacity of generation mix by carrier across Scenarios A, B, and C.}
    \label{fig:sect_coupl_capacity_all_scenarios}
\end{figure}



\subsection{Spatial capacity expansion allocation}
The optimizer allocates hydrogen production near high-irradiance zones and coastal export hubs.

Figure~\ref{fig:h2_investment_maps} shows the spatial allocation of hydrogen production capacity and supporting infrastructure across the three scenarios, highlighting the concentration of investments near high-irradiance areas and coastal export hubs.

\begin{figure}[htbp]
     \centering
     \begin{subfigure}[b]{0.3\textwidth}
         \centering
         \includegraphics[width=\textwidth]{Figures/Sector_Couple_H2/fig_capacity_map_scenario_a.pdf}
         \caption{Scenario A}
     \end{subfigure}
     \hfill
     \begin{subfigure}[b]{0.3\textwidth}
         \centering
         \includegraphics[width=\textwidth]{Figures/Sector_Couple_H2/fig_capacity_map_scenario_b.pdf}
         \caption{Scenario B}
     \end{subfigure}
     \hfill
     \begin{subfigure}[b]{0.3\textwidth}
         \centering
         \includegraphics[width=\textwidth]{Figures/Sector_Couple_H2/fig_capacity_map_scenario_b.pdf}
         \caption{Scenario C}
     \end{subfigure}
        \caption{Spatial capacity allocation for hydrogen production and supporting infrastructure.}
        \label{fig:h2_investment_maps}
\end{figure}

\subsection{System cost breakdown}
The total system expenditure exhibits notable variations across the evaluated scenarios, driven by the structural allocation of capital assets and operational expenses. To isolate the financial impact of large-scale hydrogen integration, annual expenditures are evaluated under both active export conditions and zero-export baseline conditions.

\begin{figure}[htbp]
\centering
\begin{subfigure}[b]{0.49\textwidth}
    \centering
    \includegraphics[width=\textwidth]{Figures/Sector_Couple_H2/fig_cost_by_scenario.pdf}
    \caption{With hydrogen export (70~TWh)}
    \label{fig:inv_costs_H2_eport}
\end{subfigure}
\hfill
\begin{subfigure}[b]{0.49\textwidth}
    \centering
    \includegraphics[width=\textwidth]{Figures/Sector_Couple_H2/fig_No_H2_export_cost_by_scenario.pdf}
    \caption{Zero-export baseline case}
    \label{fig:inv_costs_base}
\end{subfigure}
\caption{Total annual system cost comparison across the evaluated scenarios: (a) with hydrogen export at 70~TWh and (b) zero-export baseline case.}
\label{fig:annual_system_costs}
\end{figure}

Under the 70~TWh export demand (Fig.~\ref{fig:inv_costs_H2_eport}), total annual system expenditures range from 4.56 to 10.34~billion~\texteuro/yr. Scenario~A is heavily CAPEX-driven, dominated by dedicated renewable buildout (4.11~billion~\texteuro/yr Onshore Wind, 1.53~billion~\texteuro/yr Solar PV) and electrolysis infrastructure (3.39~billion~\texteuro/yr). In contrast, Scenarios~B and C shift financial burden toward operational expenditures (2.80 and 2.84~billion~\texteuro/yr, respectively), substituting massive wind/PV investments with grid-supplied electricity and Steam Methane Reforming with Carbon Capture (0.93~billion~\texteuro/yr in Scenario~B). Conversely, the zero-export baseline (Fig.~\ref{fig:inv_costs_base}) exhibits substantially lower overall system costs (1.44 to 2.35~billion~\texteuro/yr), where expenditures across all scenarios remain primarily tied to satisfying domestic electricity demand via modest renewable capacity, legacy oil generation, and minimal operational overhead.

\begin{figure}[htbp]
\centering
\begin{subfigure}[b]{0.49\textwidth}
    \centering
    \includegraphics[width=\textwidth]{Figures/Sector_Couple_H2/fig_opex_cost_by_scenario.pdf}
    \caption{With hydrogen export (70~TWh)}
    \label{fig:opex_costs_H2_export}
\end{subfigure}
\hfill
\begin{subfigure}[b]{0.49\textwidth}
    \centering
    \includegraphics[width=\textwidth]{Figures/Sector_Couple_H2/fig_No_H2_export_opex_cost_by_scenario.pdf}
    \caption{Zero-export baseline case}
    \label{fig:opex_costs_base}
\end{subfigure}
\caption{Operational Expenditure (OPEX) breakdown across the evaluated scenarios: (a) with hydrogen export at 70~TWh and (b) zero-export baseline case.}
\label{fig:annual_opex_costs}
\end{figure}

The OPEX breakdown (Fig.~\ref{fig:annual_opex_costs}) highlights a stark divergence between export and domestic-only cases. Under 70~TWh export demand (Fig.~\ref{fig:opex_costs_H2_export}), total OPEX surges from 0.30~billion~\texteuro/yr in Scenario~A to nearly 2.84~billion~\texteuro/yr in Scenarios~B and C, overwhelmingly driven by natural gas costs for SMR and thermal generation. Conversely, without exports (Fig.~\ref{fig:opex_costs_base}), system-wide OPEX remains below 0.60~billion~\texteuro/yr across all scenarios, driven predominantly by baseline legacy fuel demand. Notably, across all cases, primary fuel supply (gas and oil) overwhelmingly dominates operational expenditure, rendering non-fuel OPEX components negligible.
\subsection{Levelized cost of hydrogen (LCOH)}
To evaluate the economic viability of hydrogen production for export, we compute the net levelized cost of hydrogen across the considered scenarios. This metric provides a standardized basis for comparing production economics under varying grid integration strategies.

To isolate the incremental cost of hydrogen production from existing domestic power grid requirements, we utilize a paired-run subtraction method. For each scenario, the model is solved under the 70~TWh export demand (Fig.~\ref{fig:inv_costs_H2_eport}) and compared against a corresponding zero-export baseline (Fig.~\ref{fig:inv_costs_base}). The Net Levelized Cost of Hydrogen (Net LCOH) is defined as:
\begin{equation}
\label{Eq:Net-LCOH}
    \text{Net LCOH} = \frac{\text{Total Cost}_{\text{Export}} - \text{Total Cost}_{\text{Baseline}}}{\text{Exported Hydrogen Energy}}
\end{equation}
This approach ensures that domestic electricity consumers are not modeled as subsidizing the export infrastructure, providing a true marginal cost of hydrogen for the European market.


\begin{table}[htbp]
\centering
\caption{Summary of $\text{H}_2$ Export Costs and Net LCOH Breakdown}
\label{tab:h2_export_costs}
\begin{tabular}{lccc}
\toprule
Metric & Scenario A & Scenario B & Scenario C \\
\midrule
Export Run Cost [Billion \texteuro] & 10.342 & 5.923 & 4.555 \\
Baseline Zero-Export Cost [Billion \texteuro] & 2.347 & 2.308 & 1.438 \\
Net $\text{H}_2$ Cost [Billion \texteuro] & 7.995 & 3.615 & 3.117 \\
\midrule
Net LCOH [\texteuro/MWh] & 114.21 & 51.64 & 44.53 \\
Net LCOH LHV [\texteuro/kg] & 3.81 & 1.72 & 1.48 \\
Net LCOH HHV [\texteuro/kg] & 4.50 & 2.04 & 1.75 \\
\bottomrule
\multicolumn{4}{p{0.88\textwidth}}{\footnotesize \textit{Note:} Net $\text{H}_2$ Cost = Export System Cost -- Baseline Zero-Export Cost; Net LCOH = Net $\text{H}_2$ Cost / Exported $\text{H}_2$ Energy (70 TWh). Conversions assume $\text{H}_2$ Lower Heating Value (LHV) = $0.03333~\text{MWh/kg}$ and Higher Heating Value (HHV) = $0.03941~\text{MWh/kg}$.}
\end{tabular}
\end{table}

As summarized in Table~\ref{tab:h2_export_costs}, the Net LCOH exhibits significant divergence across the evaluated configurations. Across the 70~TWh export target, Scenario A yields the highest Net LCOH at 114.21~\texteuro/MWh (3.81~\texteuro/kg LHV), driven by the requirement for full renewable off-grid system expansion and dedicated asset over-provisioning. In contrast, Scenario B substantially reduces this metric to 51.64~\texteuro/MWh (1.72~\texteuro/kg LHV) through grid interaction and shared flexibility. Scenario C achieves the lowest overall production cost at 44.53~\texteuro/MWh (1.48~\texteuro/kg LHV), highlighting the strong economic advantage of co-locating generation and leveraging existing national transmission infrastructure.

\subsection{Net present value (NPV)}
Table~\ref{tab:npv_summary} presents the financial performance under a dynamic margin framework, applying a constant revenue premium of +€1.00/kg above the levelized cost of hydrogen (LCOH LHV) for each scenario over the 21-year operational horizon (2030–2050 at $r = 7.10\%$).
\begin{table}[ht]
\centering
\caption{Annual operational cash flows, project return metrics, and multi-year discounted financial evaluation by scenario (2030–2050) under a dynamic cost-plus revenue margin pricing model ($+\text{€}1.00/\text{kg}$ over scenario-specific LCOH at LHV) and an assumed zero $\text{CO}_2$ emission tax, evaluated at a discount rate of $r = 7.10\%$.}
\label{tab:npv_summary}
\small
\setlength{\tabcolsep}{8pt}
\resizebox{\linewidth}{!}{%
\begin{tabular}{lccc}
\hline
{Financial Metric} & {Scenario A} & {Scenario B} & {Scenario C} \\ \hline
LCOH LHV [EUR/kg] & 3.81 & 1.72 & {1.48} \\
Selling Price [EUR/kg] & 4.81 & 2.72 & {2.48} \\
Annual Cost [B EUR/yr] & 10.34 & 5.92 & {4.56} \\
Annual Revenue [B EUR/yr] & 10.09 & 5.72 & {5.22} \\
Annual Net Cash Flow [B EUR/yr] & -0.25 & -0.21 & {+0.66} \\
Discounted Cost [B EUR] & 111.16 & 63.66 & {48.97} \\
Discounted Revenue [B EUR] & 108.51 & 61.43 & {56.08} \\
Total System NPV [B EUR] & -2.65 & -2.23 & {+7.11} \\
Project IRR (\%) & 10.23\% & 26.30\% & {32.23\%} \\
Payback Period [Years] & 8.51 & 3.77 & {3.09} \\ \hline
\end{tabular}%
}
\end{table}

Under this pricing strategy, Scenario C demonstrates clear economic superiority by yielding a strong positive total system net present value (NPV). Achieving the lowest levelized cost of hydrogen at lower heating value (LCOH LHV) of €1.48/kg and a selling price of €2.48/kg, Scenario C delivers an annual revenue of €5.22~billion/yr against an annual system cost of €4.56~billion/yr, resulting in a positive annual net cash flow of +€0.66~billion/yr and a total system NPV of +€7.11~billion. Furthermore, Scenario C achieves an outstanding Project IRR of 32.23\% and rapid capital recovery with a payback period of just 3.09 years, significantly outperforming Scenario B (26.30\% IRR; 3.77-year payback) and Scenario A (10.23\% IRR; 8.51-year payback).

By contrast, Scenarios A and B yield negative annual net cash flows (-€0.25~billion/yr and -€0.21~billion/yr) and negative total system NPVs (-€2.65~billion and -€2.23~billion, respectively) under the +€1.00/kg margin assumption due to their higher discounted system costs (€111.16~billion and €63.66~billion). Selling at €2.72/kg in Scenario B and €4.81/kg in Scenario A is insufficient to cover total system-wide expenditures over the 21-year period. Consequently, Scenario C remains the most financially resilient option, generating the strongest returns while keeping market-clearing hydrogen export prices competitive.

\subsection{Environmental evaluation}
In this section, the performance of the model is evaluated in terms of the environmental factors.

\subsubsection{CO$_2$ emissions}
\label{sec:co2_emis}
As detailed in Table~\ref{tab:co2_mass_balance}, gross $\text{CO}_2$ generation is partitioned into net atmospheric emissions (vented) and geologically sequestered emissions (captured via CCS). Scenarios A and B operate well within the carbon cap ceiling, resulting in identical net atmospheric emissions of 6.82~million tonnes (Mt) while relying on CCUS infrastructure to manage total system generation (capturing 25.25~Mt and 43.56~Mt, respectively). Conversely, the unconstrained baseline in Scenario C yields a significantly higher total gross generation of 135.07~Mt, consisting of 59.15~Mt in atmospheric venting and 75.91~Mt in carbon capture.

\begin{table}[htbp]
\centering
\caption{Gross $\text{CO}_2$ generation partitioned into net atmospheric emissions (vented) and geologically sequestered emissions (captured via CCS), reported in million tonnes (Mt).}
\label{tab:co2_mass_balance}
\begin{tabular}{lrrr}
\toprule
{$\text{CO}_2$ Accounting Metric} & {Scenario A} & {Scenario B} & {Scenario C} \\ \midrule
Net Atmospheric $\text{CO}_2$ Emissions (Vented) & 6.82 Mt & 6.82 Mt & 59.15 Mt \\
Geologically Sequestered $\text{CO}_2$ (Captured) & 25.25 Mt & 43.56 Mt & 75.91 Mt \\ \midrule
Total Gross $\text{CO}_2$ Generated (Vented + Stored) & 32.07 Mt & 50.39 Mt & 135.07 Mt \\ \bottomrule
\end{tabular}
\end{table}


\subsubsection{Water consumption}
\label{sec:desalination}
A critical factor in the feasibility of large-scale hydrogen production in arid regions like Libya is the required water supply. In this study, the costs and energy overhead associated with Reverse Osmosis (RO) desalination are integrated directly into the hydrogen supply chain assumptions, accounting for the technology-specific water intensities of electrolysis, SMR with carbon capture, and conventional SMR.

Our analysis indicates that across all three scenarios, total hydrogen output remains virtually identical at approximately $2.13~\text{Mt}$ ($70.91\text{--}70.92~\text{TWh}_{\text{H2}}$). However, water demand varies significantly based on the technology mix deployed, requiring $21.27~\text{Mm}^3$ in Scenario A, $15.96~\text{Mm}^3$ in Scenario B, and $10.65~\text{Mm}^3$ in Scenario C (Table~\ref{tab:water_demand}).

\begin{table}[htbp]
\centering
\caption{Estimated water demand and desalination costs across scenarios.}
\label{tab:water_demand}
\begin{tabular}{lccc}
\toprule
Metric & Scenario A & Scenario B & Scenario C \\
\midrule
$\text{H}_2$ Produced ($\text{TWh}_{\text{H2}}$) & 70.91 & 70.91 & 70.92 \\
$\text{H}_2$ Mass (Million Tonnes) & 2.127 & 2.128 & 2.128 \\
Water Demand (Million $\text{m}^3$) & 21.27 & 15.96 & 10.65 \\
Est. Desalination Cost (\text{M\texteuro/yr}) & 31.91 & 23.94 & 15.97 \\
\bottomrule
\end{tabular}
\end{table}

Despite the scale of these volumes, the economic impact of water supply is remarkably low. The estimated desalination costs contribute between 1.04\% and 1.36\% to the total annualized system cost, representing a marginal addition of only \texteuro 0.01--0.02 per kg of hydrogen (€0.0075--0.0150/kg). 

From an energy perspective, the overhead required for RO desalination results in a specific energy consumption of just $0.02\text{--}0.04~\text{kWh/kg}_{\text{H2}}$ ($0.0175\text{--}0.0350~\text{kWh/kg}_{\text{H2}}$). When evaluated against total system electricity demand, desalination accounts for merely $0.03\%\text{--}0.07\%$ ($0.033\%\text{--}0.067\%$). These findings confirm that while water sourcing requires geographical planning (e.g., coastal co-location of RO plants), water consumption does not act as an economic or energetic bottleneck for Libya's hydrogen supply potential.

\subsection{System cost dynamics under $\text{CO}_2$ taxation}
The inclusion of a $100\text{~\texteuro/tonne}$ $\text{CO}_2$ tax fundamentally alters system cost dynamics, shifting economic penalties away from capital expenditure and toward carbon-intensive operation.

The impact of carbon pricing on total system costs across export scenarios is illustrated in Figure~\ref{fig:cost_breakdown_CO2_Tax}.

\begin{figure}[htbp]
    \centering
    \begin{subfigure}[b]{0.49\textwidth}
        \centering
        \includegraphics[width=\textwidth]{Figures/Sector_Couple_H2/fig_cost_by_scenario.pdf}
        \caption{Without carbon tax}
        \label{fig:cost_no_tax}
    \end{subfigure}
    \hfill
    \begin{subfigure}[b]{0.49\textwidth}
        \centering
        \includegraphics[width=\textwidth]{Figures/Sector_Couple_H2_CO2Tax100/fig_cost_by_scenario.pdf}
        \caption{With $100\text{~\texteuro/tonne}$ $\text{CO}_2$ tax}
        \label{fig:cost_with_tax}
    \end{subfigure}
    \caption{Annualized total system cost breakdown across hydrogen export scenarios ($70\text{~TWh}$) under (a) untaxed conditions and (b) a $100\text{~\texteuro/tonne}$ $\text{CO}_2$ carbon tax.}
    \label{fig:cost_breakdown_CO2_Tax}
\end{figure}

\begin{itemize}
    \item {Tax Liability vs. CAPEX Expansion:} Unabated fossil reliance drives prohibitive penalty costs under carbon pricing. In Scenario~C, the $\text{CO}_2$ tax alone reaches $5.92\text{~Billion~\texteuro/year}$ ($5.9153\text{~Billion~\texteuro/year}$), constituting over $56\%$ ($56.49\%$) of its total annual expenditure ($10.47\text{~Billion \texteuro/year}$). Conversely, higher capital investments (\text{CAPEX}) in renewables and electrolysis (Scenario~A: $1.52\text{~Billion \texteuro/year}$ Solar PV, $4.11\text{~Billion~\texteuro/year}$ Wind, $3.39\text{~Billion \texteuro/year}$ Electrolysis) limit carbon tax liabilities to just $0.68\text{~Billion \texteuro/year}$, increasing total system expenditure modestly from $10.34\text{~Billion \texteuro/year}$ to $11.02\text{~Billion \texteuro/year}$.
    \item {Technology Competitiveness:} The $100\text{~\texteuro/tonne}$ threshold erodes the cost advantage of unabated Steam Methane Reforming (\text{SMR}). Scenario~B demonstrates that integrating Carbon Capture (\text{SMR-CC}) yields the lowest overall system cost ($6.61\text{~Billion \texteuro/year}$), balancing modest capital investment ($0.93\text{~Billion \texteuro/year}$) with a heavily reduced carbon tax penalty ($0.68\text{~Billion \texteuro/year}$).
    \item {Expenditure Composition:} Under carbon pricing, high-carbon pathways (Scenario~C) transition to an operational-liability model dominated by tax payments, increasing total costs from $4.56\text{~Billion \texteuro/year}$ to $10.47\text{~Billion \texteuro/year}$, whereas decarbonized options shift system value toward fixed, asset-driven capital expenditure.
\end{itemize}

\subsection{Financial viability under carbon taxation}
A comparative summary of the core financial metrics under untaxed conditions ($0\text{~EUR/tCO}_2$) versus a $100\text{~EUR/tonne}$ $\text{CO}_2$ tax regime is presented in Table~\ref{tab:co2_tax_financial_comparison}.

\begin{table}[ht]
\centering
\small
\resizebox{\linewidth}{!}{%
\begin{threeparttable}
\caption{Comparison of key financial performance metrics across scenarios under zero tax versus a $100\text{~EUR/tonne}$ $\text{CO}_2$ tax ($r = 7.10\%$, 2030--2050 horizon, revenue margin $+1\text{~EUR/kg}$).}
\label{tab:co2_tax_financial_comparison}
\setlength{\tabcolsep}{6pt}
\begin{tabular}{lcccccc}
\toprule
& \multicolumn{2}{c}{Scenario A} & \multicolumn{2}{c}{Scenario B} & \multicolumn{2}{c}{Scenario C} \\ 
\cmidrule(lr){2-3} \cmidrule(lr){4-5} \cmidrule(lr){6-7}
Financial Metric & No Tax & Taxed & No Tax & Taxed & No Tax & Taxed \\ 
\midrule
Annual System Cost [B EUR/yr] & 10.34 & 11.02 & 5.92 & 6.61 & 4.56 & 10.47 \\
Total System NPV [B EUR] & -2.65 & -9.99 & -2.23 & -9.56 & +7.11 & -56.47 \\
Project IRR (\%) & 10.23\% & 9.24\% & 26.30\% & 20.49\% & 32.23\% & N/A\tnote{*} \\
Simple Payback Period [Years]\tnote{a} & 8.51 & 9.13 & 3.77 & 4.78 & 3.09 & N/A\tnote{*} \\ 
Discounted Payback Period [Years]\tnote{b} & 13.52 & 15.23 & 4.55 & 6.05 & 3.62 & N/A\tnote{*} \\ 
\bottomrule
\end{tabular}

\begin{tablenotes}
\footnotesize
\footnotesize
\item[a] Simple Payback Period represents the undiscounted capital recovery timeframe in years, calculated as $\text{Simple Payback} = \frac{\text{Net Overnight CAPEX}}{\text{Annual Net Operating Cash Flow}}$.
\item[b] Discounted Payback Period (DPP) accounts for the time value of money ($r = 7.10\%$), evaluated analytically as $\text{DPP} = \frac{-\ln\left(1 - \frac{\text{CAPEX} \cdot r}{\text{Annual Net Cash Flow}}\right)}{\ln(1 + r)}$.
\item[*] N/A indicates that net annual operating cash flows are strictly negative (or annual discounted inflows are insufficient to recover capital costs at $r = 7.10\%$), rendering payback unattainable and IRR non-viable.
\end{tablenotes}
\end{threeparttable}%
}
\end{table}

The introduction of carbon pricing severely undermines project bankability across all pathways, though to vastly different degrees:

\begin{itemize}
    \item {Catastrophic Impact on Fossil Pathways:} Scenario~C experiences total economic collapse under carbon pricing. The $5.92\text{~Billion EUR/year}$ ($5.915\text{~B EUR/yr}$) tax liability drives total annual system costs up from $4.56\text{~Billion EUR/year}$ to $10.47\text{~Billion EUR/year}$, rendering the total system NPV deeply negative at $-56.47\text{~Billion EUR}$ (down from $+7.11\text{~Billion EUR}$). Because annual net operating cash flows fall to $-3.40\text{~Billion EUR/year}$, the system never reaches a financial breakeven point, yielding non-viable (N/A) figures for IRR and Payback Period.
    \item {Resilience of Decarbonized Pathways:} Scenarios~A and B absorb the tax policy far more effectively due to minimal residual emissions. While both experience total system NPV contractions (Scenario~A moving from $-2.65\text{~B EUR}$ to $-9.99\text{~B EUR}$; Scenario~B from $-2.23\text{~B EUR}$ to $-9.56\text{~B EUR}$), they preserve strong asset performance and positive net operational cash flows. Scenario~B retains an exceptional financial profile with a $20.49\%$ Project IRR ($20.49\%$) and a rapid payback period of $4.78\text{~years}$.
\end{itemize}

% =====================================================================
\section{Results: Hybrid pathway for hydrogen export}
\label{sec:results_hybrid_h2} 
This section examines the results obtained under the hybrid scenario, which combines a 70~TWh hydrogen export target with a $100\text{~\texteuro/tonne}$ $\text{CO}_2$ carbon tax. The analysis focuses on the resulting generation capacity, spatial capacity allocation, and system costs under this combined policy and export setting.

\subsection{Energy mix by carrier}
The cross-scenario comparison highlights how power supply dispatch transitions across technology carriers under varying system constraints.

Figure~\ref{fig:bars_capacity_hybrid} presents the generation capacity mix by carrier under the hybrid scenario, illustrating the contribution of different technology carriers to the overall energy supply.

\begin{figure}[ht]
    \centering
    \includegraphics[width=0.8\textwidth]{Figures/Sector_Couple_H2_Hybrid/fig_Hybrid_physical_generation_by_carrier_filtered}
    \caption{Capacity of generation mix by carrier under hybrid scenario}
    \label{fig:bars_capacity_hybrid}
\end{figure}

\subsection{Spatial capacity expansion}
The optimizer allocates hydrogen production near high-irradiance zones and coastal export hubs.

Figure~\ref{fig:map_capacity_hybrid} depicts the geographical distribution of hydrogen production capacity and associated infrastructure under the hybrid scenario, revealing the spatial pattern of investment across the system.

\begin{figure}[ht]
    \centering
    \includegraphics[width=0.8\textwidth]{Figures/Sector_Couple_H2_Hybrid/fig_capacity_map_hybrid.pdf}
    \caption{Spatial capacity allocation for hydrogen production and supporting infrastructure under the hybrid scenario}
    \label{fig:map_capacity_hybrid}
\end{figure}


\subsection{System cost breakdown}
Integrating a combined policy framework—pairing a $70\text{~TWh}$ hydrogen export mandate with a $100\text{~\texteuro/tonne}$ $\text{CO}_2$ carbon tax—reshapes the overall system economics by striking a balance between capital expenditure (\text{CAPEX}) investments and operational emission penalties.

The annualized system cost structure for the hybrid scenario is shown in Figure~\ref{fig:cost_breakdown_hybrid}, comparing the $70\text{~TWh}$ export case against the zero-export baseline.

\begin{figure}[htbp]
    \centering
    \begin{subfigure}[b]{0.40\textwidth}
        \centering
        \includegraphics[width=\textwidth]{Figures/Sector_Couple_H2_Hybrid/fig_cost_by_scenario.pdf}
        \caption{With hydrogen export (70~TWh)}
        \label{fig:cost_hybrid_export}
    \end{subfigure}
    \hfill
    \begin{subfigure}[b]{0.40\textwidth}
        \centering
        \includegraphics[width=\textwidth]{Figures/Sector_Couple_H2_Hybrid/fig_No_H2_export_cost_by_scenario}
        \caption{Zero-export baseline case}
        \label{fig:cost_hybrid_no_export}
    \end{subfigure}
\caption{Annualized total system cost breakdown under the hybrid scenario ($100\text{~\texteuro/tonne}$ $\text{CO}_2$ carbon tax): (a) with 70~TWh hydrogen export and (b) zero-export baseline case.}
\label{fig:cost_breakdown_hybrid}
\end{figure}

\begin{itemize}
    \item {Export Burden and Scale Expansion:} Meeting the $70\text{~TWh}$ external hydrogen demand elevates the total annual system cost to $7.47\text{~Billion~\texteuro/year}$ ($7.469\text{~B EUR/yr}$), representing a $110.75\%$ increase over the zero-export baseline ($3.54\text{~Billion~\texteuro/year}$). This expenditure surge is overwhelmingly driven by natural gas feedstock procurement and operational scaling rather than pure renewable expansion.
    \item {SMR with Carbon Capture Dominance:} To satisfy the export volume under strict carbon penalties without over-investing in dedicated renewable generation, the optimizer relies heavily on Steam Methane Reforming with Carbon Capture (\text{SMR-CC}). Capital allocation for \text{SMR-CC} stands as the largest single technology investment at $0.93\text{~Billion~\texteuro/year}$ ($0.9280\text{~B EUR/yr}$ in export vs.~$0.01\text{~B EUR/yr}$ baseline), whereas green electrolysis CAPEX remains negligible ($0.0003\text{~Billion~\texteuro/year}$).
    \item {Operational Cost and Tax Liabilities:} Operating costs constitute the single largest cost component in the export case at $3.60\text{~Billion~\texteuro/year}$ ($48.26\%$ of total system expenditure), reflecting fuel inputs for fossil-based pathways. Meanwhile, the $100\text{~\texteuro/tonne}$ carbon tax liability increases from $0.59\text{~Billion~\texteuro/year}$ in the zero-export baseline to $0.79\text{~Billion~\texteuro/year}$ in the export case, demonstrating that carbon capture deployment effectively caps emission penalties despite the massively scaled hydrogen supply output.
    \item {Renewable Integration Profile:} Onshore wind forms the backbone of electricity CAPEX ($1.15\text{~Billion~\texteuro/year}$ under export vs.~$1.15\text{~Billion~\texteuro/year}$ baseline), outscaling solar PV investment ($0.33\text{~Billion~\texteuro/year}$ across both cases) in both operational regimes.
\end{itemize}

\begin{figure}[htbp]
\centering
\begin{subfigure}[b]{0.40\textwidth}
    \centering
    \includegraphics[width=\textwidth]{Figures/Sector_Couple_H2_Hybrid/fig_opex_cost_hybrid.pdf}
    \caption{With hydrogen export (70~TWh)}
    \label{fig:hybrid_opex_costs_H2_export}
\end{subfigure}
\hfill
\begin{subfigure}[b]{0.40\textwidth}
    \centering
    \includegraphics[width=\textwidth]{Figures/Sector_Couple_H2_Hybrid/fig_No_H2_export_opex_cost_hybrid.pdf}
    \caption{Zero-export baseline case}
    \label{fig:hybrid_opex_costs_base}
\end{subfigure}
\caption{Operational Expenditure (OPEX) breakdown for the hybrid model ($100\text{~\texteuro/tonne}$ $\text{CO}_2$ carbon tax): (a) with 70~TWh hydrogen export and (b) zero-export baseline case.}
\label{fig:hybrid_annual_opex_costs}
\end{figure}

The hybrid OPEX breakdown (Fig.~\ref{fig:hybrid_annual_opex_costs}) illustrates that 70~TWh export demand surges operational expenses from $1.50\text{~Billion~\texteuro/year}$ to $4.40\text{~Billion~\texteuro/year}$, heavily dominated by natural gas input for SMR with CC ($2.78\text{~Billion~\texteuro/year}$ vs.~$0.28\text{~Billion~\texteuro/year}$ baseline), alongside consistent oil fuel costs ($0.62\text{~Billion~\texteuro/year}$) and carbon tax liabilities ($0.79\text{~Billion~\texteuro/year}$ export vs.~$0.59\text{~Billion~\texteuro/year}$ baseline). Notably, primary fuel procurement and carbon penalties overwhelmingly account for the system's operational expenditure, rendering non-fuel OPEX components negligible across both export regimes.


\subsection{Levelized cost of hydrogen (LCOH) under hybrid pathway}
To evaluate the economic viability of hydrogen production for export under combined carbon pricing and export constraints, we compute the net levelized cost of hydrogen for the hybrid scenario. This metric provides a standardized basis for evaluating export competitiveness under a $100\text{~\texteuro/tonne}$ $\text{CO}_2$ carbon tax regime.

To isolate the incremental cost of hydrogen production from existing domestic power grid requirements, we utilize a paired-run subtraction method based on Eq.~\ref{Eq:Net-LCOH}. The model is solved under the 70~TWh export demand and compared against a corresponding zero-export baseline.

This subtraction approach ensures that domestic electricity consumers are not modeled as subsidizing the export infrastructure, providing a true marginal production cost for the European market.

\textit{Remark:} The LCOH calculations presented here reflect purely capital and operational production costs and explicitly exclude $\text{CO}_2$ carbon tax liabilities.

\begin{table}[htbp]
\centering
\caption{Summary of $\text{H}_2$ Export costs and net LCOH breakdown for the hybrid pathway}
\label{tab:h2_export_costs_hybrid}
\begin{tabular}{lc}
\toprule
Metric & Hybrid Scenario \\
\midrule
Export Run Cost [Billion \texteuro] & 6.675 \\
Baseline Zero-Export Cost [Billion \texteuro] & 2.949 \\
Net $\text{H}_2$ Cost [Billion \texteuro] & 3.725 \\
Exported $\text{H}_2$ Energy [TWh] & 70.0 \\
\midrule
Net LCOH [\texteuro/MWh] & 53.22 \\
Net LCOH LHV [\texteuro/kg] & 1.77 \\
Net LCOH HHV [\texteuro/kg] & 2.10 \\
\bottomrule
\multicolumn{2}{p{0.65\textwidth}}{\scriptsize \textit{Note:} Net $\text{H}_2$ Cost = Export System Cost -- Baseline Zero-Export Cost; Net LCOH = Net $\text{H}_2$ Cost / Exported $\text{H}_2$ Energy. LCOH figures exclude $\text{CO}_2$ carbon tax liabilities. Conversions assume $\text{H}_2$ Lower Heating Value (LHV) = $0.03333~\text{MWh/kg}$ and Higher Heating Value (HHV) = $0.03941~\text{MWh/kg}$.}
\end{tabular}
\end{table}

As summarized in Table~\ref{tab:h2_export_costs_hybrid}, the hybrid policy environment results in a Net LCOH of $53.22\text{~\texteuro/MWh}$, which translates to $1.77\text{~\texteuro/kg}$ on a Lower Heating Value (LHV) basis and $2.10\text{~\texteuro/kg}$ on a Higher Heating Value (HHV) basis. Subtracting the baseline domestic grid cost ($2.95\text{~Billion~\texteuro}$) from the full export system expenditure ($6.67\text{~Billion~\texteuro}$) yields an incremental export system cost of $3.73\text{~Billion~\texteuro}$ ($3.725\text{~Billion~\texteuro}$).

This unit production cost highlights the economic trade-off of the hybrid regime: while the $100\text{~\texteuro/tonne}$ $\text{CO}_2$ tax restricts unabated fossil options and forces investment in carbon capture infrastructure (\text{SMR-CC}), the ability to share grid flexibility prevents the extreme capital premiums associated with purely off-grid renewable pathways. Consequently, the resulting price point of $1.77\text{~\texteuro/kg LHV}$ positions Libyan hydrogen exports competitively against projected European import cost benchmarks.

\subsection{Net present value (NPV) under hybrid pathway}
Table~\ref{tab:npv_summary_hybrid} presents the financial performance metrics for the hybrid scenario under a dynamic cost-plus pricing framework. A constant revenue margin of $+1.00\text{~\texteuro/kg}$ above the scenario LCOH LHV is applied over the 21-year operational horizon (2030–2050) using a $7.10\%$ real discount rate. While the base LCOH metric reflects unit production costs exclusive of carbon penalties, carbon tax liabilities are separately integrated into the broader system-level cash flow, NPV, IRR, and payback evaluations to ensure comprehensive financial accounting.


\begin{table}[htbp]
\centering
\begin{threeparttable}
\caption{Annual operational cash flows, project return metrics, and multi-year discounted financial evaluation (2030--2050) for the hybrid scenario under a $+1.00\text{~\texteuro/kg}$ revenue margin pricing model and a $100\text{~\texteuro/tonne}$ $\text{CO}_2$ carbon tax ($r = 7.10\%$).}
\label{tab:npv_summary_hybrid}
\small
\setlength{\tabcolsep}{12pt}
\begin{tabular}{lc}
\toprule
Financial Metric & Hybrid Scenario \\ 
\midrule
LCOH LHV Cost [\texteuro/kg] & 1.77 \\
Selling Price (+1 \texteuro/kg) [\texteuro/kg] & 2.77 \\
Export System Cost [B \texteuro/yr] & 6.67 \\
$\text{CO}_2$ Tax Liability [B \texteuro/yr] & 0.79 \\
Annual Total System Cost [B \texteuro/yr] & 7.47 \\
Annual Revenue [B \texteuro/yr] & 5.83 \\
Annual System Net Cash Flow [B \texteuro/yr] & -1.64 \\
Annual Project Net Cash Inflow [B \texteuro/yr] & +1.81 \\
\midrule
Discounted System Cost [B \texteuro] & 80.29 \\
Discounted Gross Revenue [B \texteuro] & 62.62 \\
Total System NPV [B \texteuro] & -17.67 \\
Project IRR (\%) & 20.69\% \\
Simple Payback Period [Years]\tnote{a} & 4.74 \\ 
Discounted Payback Period [Years]\tnote{b} & 5.98 \\ 
\bottomrule
\end{tabular}

\begin{tablenotes}
\footnotesize
\item[*] Revenue margin is added directly to LCOH LHV. Annual total system cost equals Export System Cost plus $\text{CO}_2$ Tax Liability. System annual costs and subsequent cash flows explicitly include the $100\text{~\texteuro/tonne}$ $\text{CO}_2$ carbon tax liability.
\item[a] Simple Payback Period represents the undiscounted capital recovery timeframe in years, calculated as $\text{Simple Payback} = \frac{\text{Net Overnight CAPEX}}{\text{Annual Net Operating Cash Flow}}$.
\item[b] Discounted Payback Period (DPP) accounts for the time value of money ($r = 7.10\%$), evaluated analytically as $\text{DPP} = \frac{-\ln\left(1 - \frac{\text{CAPEX} \cdot r}{\text{Annual Net Cash Flow}}\right)}{\ln(1 + r)}$.
\end{tablenotes}
\end{threeparttable}
\end{table}

Under the combined policy setting, the base LCOH LHV of $1.77\text{~\texteuro/kg}$ establishes a dynamic market selling price of $2.77\text{~\texteuro/kg}$. At this price level, annual gross revenues reach $5.83\text{~Billion~\texteuro/year}$ ($5.82568\text{~B EUR/yr}$). When evaluated against the total annual system expenditure of $7.47\text{~Billion~\texteuro/year}$ (composed of $6.67\text{~Billion~\texteuro/year}$ in export system costs and $0.79\text{~Billion~\texteuro/year}$ in carbon tax liabilities), total system-level net cash flows settle at $-1.64\text{~Billion~\texteuro/year}$, yielding a multi-year discounted system NPV of $-17.67\text{~Billion~\texteuro}$ ($62.62\text{~Billion~\texteuro}$ discounted revenue vs. $80.29\text{~Billion~\texteuro}$ discounted system cost).

However, when evaluating the commercial viability of the dedicated hydrogen export infrastructure, the $+1.00\text{~\texteuro/kg}$ margin over net production costs generates a strong dedicated annual net cash inflow of $+1.81\text{~Billion~\texteuro/year}$ ($1.8088\text{~B EUR/yr}$). Assessing this project-level cash flow stream against net initial capital outlays yields a Project Internal Rate of Return (IRR) of $20.69\%$ and a rapid payback period of $4.74\text{~years}$.

These results confirm that even under a penalizing $100\text{~\texteuro/tonne}$ carbon tax regime, the hybrid export project remains highly bankable under a standard cost-plus commercial off-take agreement. The $+1.00\text{~\texteuro/kg}$ premium provides more than sufficient financial cushion to secure strong double-digit project returns well above the $7.10\%$ hurdle rate.

\subsection{Environmental and resource evaluation under hybrid pathway}
\label{sec:environmental_eval}
This section evaluates the environmental impact and resource demand of the hybrid scenario, focusing on $\text{CO}_2$ emissions balance, carbon capture effectiveness, and fresh water requirements for hydrogen production.

Under the hybrid regime, gross system $\text{CO}_2$ generation is partitioned into net atmospheric emissions (vented) and geologically sequestered emissions (captured via CCS). As detailed in Table~\ref{tab:co2_mass_balance_hybrid}, total gross generation reaches $52.62\text{~Mt}$, of which $44.68\text{~Mt}$ ($84.91\%$) is successfully captured and sequestered. The remaining $7.94\text{~Mt}$ of net atmospheric emissions consists primarily of unabated oil generation ($5.72\text{~Mt}$), uncaptured process emissions from \text{SMR-CC} ($2.03\text{~Mt}$), and residual industrial gas usage ($0.19\text{~Mt}$). This high capture rate highlights how the $100\text{~\texteuro/tonne}$ carbon tax effectively incentivizes deep carbon abatement across the production chain.

\begin{table}[htbp]
\centering
\caption{$\text{CO}_2$ mass balance and vented emissions breakdown for the hybrid scenario.}
\label{tab:co2_mass_balance_hybrid}
\begin{tabular}{lc}
\toprule
$\text{CO}_2$ Accounting Metric / Source & Value \\
\midrule
Net Atmospheric $\text{CO}_2$ Emissions (Vented) & 7.94 Mt \\
Geologically Sequestered $\text{CO}_2$ (Captured) & 44.68 Mt \\
{Total Gross $\text{CO}_2$ Generated (Vented + Stored)} & {52.62 Mt} \\
\midrule
\textit{Breakdown of Vented Emissions by Source:} & \\
~~Oil & 5.72 Mt \\
~~\text{SMR-CC} & 2.03 Mt \\
~~Industrial Gas & 0.19 Mt \\
\bottomrule
\end{tabular}
\end{table}

A critical consideration for hydrogen production in Libya's arid environment is the associated water consumption. Reverse Osmosis (RO) desalination energy overheads and costs are integrated directly into the system model, applying technology-specific water intensities ($10\text{~L/kg}$ for electrolysis, $7.5\text{~L/kg}$ for \text{SMR-CC}, and $5\text{~L/kg}$ for unabated \text{SMR}). 

To support a total hydrogen output of $2.13\text{~Mt}$ ($70.91\text{~TWh}$), total fresh water demand is $15.96\text{~Million~m}^3/\text{year}$ (Table~\ref{tab:water_demand_hybrid}). Despite this volume, desalination imposes a negligible economic and energetic burden on the system: annual desalination costs total $23.93\text{~Million~\texteuro/year}$, accounting for just $0.68\%$ ($0.675\%$) of total system expenditure and adding a marginal cost of only $0.01\text{~\texteuro/kg}_{\text{H}_2}$ ($0.0112\text{~\texteuro/kg}_{\text{H}_2}$). From an energy perspective, the required $55.85\text{~GWh}$ of electricity consumption yields a specific energy overhead of $0.03\text{~kWh/kg}_{\text{H}_2}$ ($0.0262\text{~kWh/kg}_{\text{H}_2}$), representing a mere $0.05\%$ of overall system electricity demand. Consequently, while coastal desalination infrastructure is essential, water supply does not act as an economic or energetic bottleneck for hydrogen export.

\begin{table}[htbp]
\centering
\caption{Dynamic multi-technology water consumption and desalination summary for the hybrid scenario.}
\label{tab:water_demand_hybrid}
\begin{tabular}{lc}
\toprule
Desalination \& Water Consumption Metric & Hybrid Scenario \\
\midrule
Total $\text{H}_2$ Output ($\text{TWh}_{\text{H2}}$ / Million Tonnes) & 70.91 / 2.13 \\
Total Water Demand ($\text{Million~m}^3$) & 15.96 \\
Estimated Desalination Cost ($\text{Million~\texteuro/yr}$) & 23.93 \\
Desalination Cost Share of Total System (\%) & 0.675\% \\
Desalination Marginal Cost ($\text{\texteuro/kg}_{\text{H2}}$) & 0.0112 \\
Desalination Energy Overhead ($\text{GWh}$) & 55.85 \\
Desalination Specific Energy ($\text{kWh/kg}_{\text{H2}}$) & 0.0262 \\
Desalination Share of Total Electricity Demand (\%) & 0.050\% \\
\bottomrule
\multicolumn{2}{p{0.68\textwidth}}{\footnotesize \textit{Note:} Based on an assumed seawater desalination cost of $1.50\text{~\texteuro/m}^3$ and technology-specific water intensities.}
\end{tabular}
\end{table}

% =====================================================================

\section{Conclusion}
PyPSA-Libya provides the first open-source, high-resolution look at Libya’s energy future. The results suggest that Libya could become a competitive exporter of green hydrogen to Europe, provided that transmission bottlenecks are addressed.

% =====================================================================

\section*{Code and data availability}
The model code is available at \url{https://github.com/mhdella/pypsa-libya}. 

% =====================================================================
% Appendices   
% =====================================================================
\appendix
\section{Optimization problem scale and computational complexity}
\label{sec:appendix_problem_scale}

\setcounter{table}{0}
\renewcommand{\thetable}{A.\arabic{table}}
\setcounter{equation}{0}
\renewcommand{\theequation}{A.\arabic{equation}}
\setcounter{figure}{0}
\renewcommand{\thefigure}{A.\arabic{figure}}

The sector-coupled optimization model for Libya is formulated as a Linear Program (LP) using \texttt{linopy} with 2,920 time snapshots (8,760 hours at a 3-hour resolution) across 12 spatial clusters. Table~\ref{tab:problem_dimensions} summarizes the structural dimensions of the linear programming matrices for the three analyzed scenarios.

\begin{table}[htbp]
\centering
\caption{Linear programming matrix dimensions across the modeled scenarios.}
\label{tab:problem_dimensions}
\begin{tabular}{lrrr}
\toprule
{Scenario} & {Decision Variables} & {Constraints} & {Non-Zero Entries} \\ \midrule
Scenario A & 2,228,414 & 4,596,589 & 9,611,971 \\
Scenario B & 2,263,466 & 4,666,680 & 9,752,143 \\
Scenario C & 2,298,518 & 4,736,772 & 9,962,395 \\ \bottomrule
\end{tabular}
\end{table}

The primary driver of the dimensional variation across scenarios is the number of enabled hydrogen production technologies. In PyPSA, each active technology instantiated across network nodes generates both capacity investment variables and operational dispatch decision variables for every time step. Scenario A restricts production exclusively to \texttt{H2 Electrolysis}, Scenario B includes \texttt{H2 Electrolysis} and \texttt{SMR CC}, and Scenario C permits all three options (\texttt{H2 Electrolysis}, \texttt{SMR}, and \texttt{SMR CC}). Excluding a technology skips component instantiation; each additional SMR technology adds exactly 35,052 decision variables ($2,920 \text{snapshots} \times 12 \text{nodes} + \text{investment variables}$) and $70,092$ constraints to the matrix.

% =====================================================================

\bibliographystyle{elsarticle-num}
\bibliography{Libya_H2}

\end{document}